\documentclass[aspectratio=169]{beamer}

\usepackage{listings}
\usepackage{xcolor}

\definecolor{codegray}{rgb}{0.5,0.5,0.5}
\definecolor{codepurple}{rgb}{0.58,0,0.82}
\definecolor{backcolour}{rgb}{0.95,0.95,0.92}

\lstdefinestyle{customc}{
  backgroundcolor=\color{backcolour},
  commentstyle=\color{codegray}\itshape,
  keywordstyle=\color{blue}\bfseries,
  numberstyle=\tiny\color{gray},
  stringstyle=\color{codepurple},
  basicstyle=\ttfamily\footnotesize,
  breakatwhitespace=false,
  breaklines=true,
  captionpos=b,
  keepspaces=true,
  numbers=left,
  numbersep=5pt,
  showspaces=false,
  showstringspaces=false,
  showtabs=false,
  tabsize=2,
  frame=single,
  language=C++
}

\lstset{style=customc}

\usepackage[spanish]{babel}
\usepackage[utf8]{inputenc}
\usepackage[T1]{fontenc}
\usepackage{lmodern}
\usepackage{amsmath, amsfonts}
\usepackage{hyperref}

\usetheme{Madrid}

\title{La STL en C++}
\author{Sebastian Mestre}
\date{}

\begin{document}

\begin{frame}
  \titlepage
\end{frame}

\section{Introducción}

\begin{frame}{¿Qué es la STL?}
  \begin{itemize}
    \item C++ trae varias \textbf{estructuras de datos} y \textbf{algoritmos} ya implementados.
    \item El objetivo es \textbf{simplificarnos la vida}: usar código probado en lugar de reinventar la rueda.
    \item Algunos contenedores importantes:
      \begin{itemize}
        \item \texttt{std::vector<T>} — arreglo dinámico.
        \item \texttt{std::set<T>} — conjunto ordenado.
        \item \texttt{std::map<K,V>} — diccionario ordenado por clave.
        \item \texttt{std::stack<T>} — pila (LIFO).
        \item \texttt{std::queue<T>} — cola (FIFO).
        \item \texttt{std::priority\_queue<T>} — cola de prioridad (max heap).
      \end{itemize}
    \item Referencia completa: \href{https://en.cppreference.com/w/cpp/container}{cppreference: containers}.
  \end{itemize}
\end{frame}

\section{Iteradores}

\begin{frame}{Interfaz de iteradores}
  \begin{itemize}
    \item Muchas estructuras de la STL exponen una \textbf{interfaz de iteradores} para recorrer sus elementos.
    \item Funciones y operaciones básicas:
      \begin{itemize}
        \item \texttt{begin(s)} — iterador al inicio de una secuencia \texttt{s}.
        \item \texttt{end(s)} — iterador al final de una secuencia \texttt{s}.
        \item \texttt{*it} — elemento apuntado por el iterador \texttt{it}.
      \end{itemize}
    \item Conceptualmente, un iterador representa una \textbf{posición} en la secuencia.
    \item En una secuencia de \(N\) elementos hay \(N+1\) posiciones (antes del primero, entre cada par, y después del último).
  \end{itemize}
\end{frame}

\begin{frame}[fragile]{Recorriendo con iteradores}
  \begin{itemize}
    \item Podemos recorrer una secuencia usando un bucle \texttt{for} con iteradores.
  \end{itemize}
  \begin{block}{Ejemplo}
  \begin{lstlisting}[style=customc]
vector<int> s = {2, 3, 5, 8, 13};

for (auto it = begin(s); it != end(s); ++it) {
  cout << *it << " ";
}
cout << "\n";
  \end{lstlisting}
  \end{block}
  \begin{itemize}
    \item Operaciones útiles sobre iteradores:
      \begin{itemize}
        \item \texttt{it1 == it2} verifica si apuntan a la misma posición.
        \item \texttt{++it} / \texttt{--it} avanzar / retroceder.
        \item \texttt{prev(it)}, \texttt{next(it)} — iterador anterior / siguiente.
      \end{itemize}
  \end{itemize}
\end{frame}

\section{Ejemplo: mínima diferencia}

\begin{frame}{Problema de ejemplo}
  \begin{block}{Enunciado}
    Nos dan un arreglo de \(N\) números distintos \(a[1], a[2], \dots, a[N]\).\\
    Encontrar la \textbf{mínima diferencia} entre dos elementos de este arreglo.
  \end{block}
  \begin{itemize}
    \item Restricciones:
      \begin{itemize}
        \item \(1 \le N \le 5 \cdot 10^5\)
        \item \(1 \le a[i] \le 10^9\)
      \end{itemize}
    \item Entrada:
      \begin{itemize}
        \item Primera línea: \(N\).
        \item Segunda línea: los \(N\) valores distintos \(a[1], \dots, a[N]\).
      \end{itemize}
    \item Salida: la mínima diferencia entre dos elementos de \(a\).
  \end{itemize}
\end{frame}

\begin{frame}[fragile]{Ejemplo de entrada/salida}
  \begin{block}{Ejemplo}
  \small
  Entrada:
  \begin{verbatim}
6
9000 5004 1000 8900 5001 1100
  \end{verbatim}
  Salida:
  \begin{verbatim}
3
  \end{verbatim}
  \end{block}
\end{frame}

\begin{frame}[fragile]{Solución naive \(O(N^2)\)}
  \begin{itemize}
    \item Idea: probar \textbf{todos los pares} \((i, j)\) y quedarnos con la mínima diferencia.
    \item Complejidad: \(O(N^2)\), demasiado lenta para \(N = 5 \cdot 10^5\).
  \end{itemize}
  \begin{block}{Código}
  \begin{lstlisting}[style=customc]
int a[500100];
int main() {
  int n; cin >> n;
  forn(i, n) cin >> a[i];

  int ans = abs(a[1] - a[0]);
  forn(i, n) forn(j, n)
      if (i != j) ans = min(ans, abs(a[i] - a[j]));
  cout << ans << "\n";
}\end{lstlisting}
  \end{block}
\end{frame}

\section{std::vector}

\begin{frame}{\texttt{std::vector}: arreglo dinámico}
  \begin{itemize}
    \item \texttt{std::vector<T>} es la estructura más usada en programación competitiva.
    \item Se comporta como un \textbf{arreglo dinámico}:
      \begin{itemize}
        \item \texttt{vector<T>(n)} crea un vector de \(n\) elementos.
        \item \texttt{a[i]} acceso al elemento en la posición \(i\).
        \item \texttt{a.size()} tamaño del vector. (todas las estructuras de la STL tienen esto)
        \item \texttt{a.push\_back(x)} agrega un elemento al final.
        \item \texttt{a.pop\_back()} elimina el último elemento.
        \item \texttt{a.back()} último elemento.
        \item \texttt{a.clear()} vacia un vector
        \item \texttt{a.empty()} responde si esta vacio
        \item \texttt{begin(a)}, \texttt{end(a)} iteradores.
      \end{itemize}
    \item Referencia: \href{https://en.cppreference.com/w/cpp/container/vector}{cppreference: vector}.
  \end{itemize}
\end{frame}

\begin{frame}[fragile]{Leyendo con \texttt{vector}}
  \begin{itemize}
    \item Con \texttt{vector} no necesitamos fijar un \(N\) máximo a mano.
  \end{itemize}
  \begin{block}{Ejemplo}
  \begin{lstlisting}[style=customc]
int n; cin >> n;

vector<int> a(n);
forn(i, n) cin >> a[i];\end{lstlisting}
  \end{block}
  \begin{itemize}
    \item El resto del código puede ser igual a la versión con arreglo estático.
  \end{itemize}
\end{frame}

\section{Usando \texttt{std::set}}

\begin{frame}[fragile]{Solución con \texttt{std::set}}
  \begin{itemize}
    \item \texttt{std::set<int>} mantiene los elementos \textbf{ordenados y sin duplicados}.
    \item Podemos ir insertando elementos y comparar con sus \textbf{vecinos} en el conjunto.
  \end{itemize}
  \begin{block}{Idea}
    Para cada \(x\) en \(a\):
    \begin{itemize}
      \item Buscamos el primer elemento \(\ge x\) y el anterior a él.
      \item Actualizamos la mejor diferencia.
    \end{itemize}
  \end{block}
\end{frame}

\begin{frame}[fragile]{\texttt{std::set}}
  \begin{itemize}
    \item \texttt{std::set<T>} es un conjunto ordenado de elementos únicos.
    \item Métodos principales:
      \begin{itemize}
        \item \texttt{s.size()} cantidad de elementos
        \item \texttt{s.empty()} ¿está vacío?
        \item \texttt{s.clear()} elimina todos los elementos
        \item \texttt{s.insert(x)} inserta el elemento \(x\)
        \item \texttt{s.erase(x)} elimina el elemento \(x\)
        \item \texttt{s.count(x)} indica si \(x\) está presente (1 si está, 0 si no)
        \item \texttt{s.find(x)} devuelve un iterador a la izquierda de \(x\) (o \texttt{end(s)} si no está)
        \item \texttt{s.lower\_bound(x)} iterador que separa los elementos \(< x\) de los \(\ge x\)
        \item \texttt{s.upper\_bound(x)} iterador que separa los elementos \(\le x\) de los \(> x\)
        \item \texttt{begin(s)}, \texttt{end(s)} iteradores al inicio y final
      \end{itemize}
    \item Todos los métodos tienen costo \(O(\log N)\).
    \item Referencia: \href{https://en.cppreference.com/w/cpp/container/set}{cppreference: set}.
  \end{itemize}
\end{frame}

\begin{frame}[fragile]{Código con \texttt{set} (\(O(N \log N)\))}
  \begin{block}{Implementación}
  \begin{lstlisting}[style=customc]
int ans = abs(a[1] - a[0]);

set<int> s;
for (int x : a) {                      // O(N) iteraciones
  if (auto it = s.upper_bound(x); it != end(s)) { // O(log N)
    ans = min(ans, *it - x);           // O(1)
  }
  if (auto it = s.lower_bound(x); it != begin(s)) { // O(log N)
    ans = min(ans, x - *prev(it));     // O(1)
  }
  s.insert(x);                         // O(log N)
}
// total: O(N log N)
  \end{lstlisting}
  \end{block}
\end{frame}

\begin{frame}[fragile]{Variación con \texttt{set}}
  \begin{itemize}
    \item También podemos construir el \texttt{set} a partir del \texttt{vector}.
    \item \textbf{Observación}: todo elemento es vecino derecho de su vecino izquierdo
  \end{itemize}
  \begin{block}{Ejemplo}
  \begin{lstlisting}[style=customc]
int ans = abs(a[1] - a[0]);

set<int> s(begin(a), end(a)); // O(N log N)
for (int x : a) {
  if (auto it = s.upper_bound(x); it != end(s)) {
    ans = min(ans, *it - x);
  }
}
// total: O(N log N)
  \end{lstlisting}
  \end{block}
\end{frame}

\section{\texttt{<algorithm>}}

\begin{frame}{La cabecera \texttt{<algorithm>}}
  \begin{itemize}
    \item Además de contenedores, C++ trae una \textbf{biblioteca de algoritmos}.
    \item Algunos muy usados:
      \begin{itemize}
        \item \texttt{sort(begin(a), end(a))} ordena el vector \texttt{a}.
        \item \texttt{reverse(begin(a), end(a))} invierte el orden.
        \item \texttt{rotate(begin(a), begin(a)+k, end(a))} rota \(k\) posiciones.
        \item \texttt{unique(begin(a), end(a))} compacta elementos únicos consecutivos.
        \item \texttt{fill(begin(a), end(a), x)} llena con un valor.
      \end{itemize}
    \item Referencia: \href{https://en.cppreference.com/w/cpp/algorithm}{cppreference: algorithm}.
  \end{itemize}
\end{frame}

\begin{frame}[fragile]{Solución óptima con \texttt{sort}}
  \begin{itemize}
    \item \textbf{Observación:} Si ordenamos el arreglo, los vecinos serán elementos consecutivos.
    \item Complejidad total: \(O(N \log N)\).
  \end{itemize}
  \begin{block}{Implementación}
  \begin{lstlisting}[style=customc]
sort(begin(a), end(a)); // O(N log N)
int ans = a[1] - a[0];
forn(i, n-1) ans = min(ans, a[i+1] - a[i]);
cout << ans << "\n";
// total: O(N log N)
  \end{lstlisting}
  \end{block}
\end{frame}

\section*{Práctica}

\begin{frame}{Problemas para practicar}
  \begin{itemize}
    \item Problemas de CSES donde la STL y \texttt{<algorithm>} son especialmente útiles:
      \begin{itemize}
        \item \href{https://cses.fi/problemset/task/1083}{Missing Number}.
        \item \href{https://cses.fi/problemset/task/1755}{Palindrome Reorder}.
        \item \href{https://cses.fi/problemset/task/2205}{Gray Code}.
        \item \href{https://cses.fi/problemset/task/1621}{Distinct Numbers}.
        \item \href{https://cses.fi/problemset/task/1091}{Concert Tickets}.
        \item \href{https://cses.fi/problemset/task/1619}{Restaurant Customers}.
      \end{itemize}
    \item Recomendación: implementar varias soluciones usando \texttt{vector}, \texttt{set}, \texttt{map} y algoritmos de la STL.
  \end{itemize}
\end{frame}

\begin{frame}{Preguntas}
  \centering
  \Huge
  \textbf{Q\&A}
\end{frame}

\end{document}